\section{模型假设}

（1）检测点坐标精确：检测点坐标视为精确已知量。

（2）干扰源静止：同一频道的历次示向度读数来自同一个静止干扰源。

（3）误差界约束：示向度误差仅由题设给定的误差界约束。

（4）扇面楔形近似：问题一忽略径向限制，将观测区域近似为楔形，并对其交集采用凸几何算法分析。

（5）三类反馈框架：检测反馈包括近距离、示向度和无信号三类，并分别更新定位区域。

（6）最小包围圆评价：以反馈后定位区域的最小包围圆圆心为预测位置、半径为误差上界。
